
\subsubsection{6.1 Interpretation of Findings}

The results support the hypothesis that RLHF instruction tuning degrades discriminative confidence quality through geometric distortion. The causal chain suggested by the evidence is:

\begin{enumerate}
\item RLHF preference optimization rewards confident-sounding responses
\item This inflates logit margins, including for incorrect predictions (H-M1)
\item Margin inflation decouples the confidence-correctness relationship (H-M2)
\item The result is degraded discriminative ability measurable as AUROC decrease (H-E1)
\item Brier decomposition confirms the distortion affects refinement, not just reliability (H-M3)
\end{enumerate}

The finding that distortion is geometric rather than scalar has implications for post-hoc calibration. Temperature scaling is designed to correct scalar miscalibration by adjusting probability magnitudes. If the distortion is geometric---affecting the shape of the confidence-correctness relationship---temperature scaling may improve ECE while leaving discriminative degradation unaddressed.

\subsubsection{6.2 Cross-Family Consistency}

The effect direction is consistent across both tested families (Qwen and Mistral), though effect magnitudes differ. Mistral shows larger margin inflation (16.79x vs. 3.06x) and greater refinement degradation (0.0487 vs. 0.0216). This magnitude asymmetry may reflect differences in RLHF training procedures, preference datasets, or base model characteristics between model developers.

The consistency of effect direction across architecturally distinct model families suggests that the phenomenon may be related to RLHF training procedures rather than vendor-specific implementation details, though this interpretation is limited by the absence of Llama family data.

\subsubsection{6.3 Limitations}

Several limitations constrain the interpretation of these results:

\begin{enumerate}
\item \textbf{Model family coverage:} The Llama-2-7B family was not tested due to HuggingFace access restrictions. Generalization to a third family remains unconfirmed.
\end{enumerate}

\begin{enumerate}
\item \textbf{Model scale:} Only 7B-parameter models were evaluated. The effect magnitude may differ at other scales.
\end{enumerate}

\begin{enumerate}
\item \textbf{Dataset scope:} Only MMLU was used. Generalization to other multiple-choice benchmarks (TruthfulQA, ARC, CommonsenseQA) and to non-multiple-choice tasks is not established.
\end{enumerate}

\begin{enumerate}
\item \textbf{Prompt format:} Only zero-shot prompting was evaluated. The planned 2x2 design (zero-shot vs. few-shot $\times$ base vs. instruct) was simplified to 1x2.
\end{enumerate}

\begin{enumerate}
\item \textbf{Temperature scaling not tested:} The claim that temperature scaling cannot repair geometric distortion is theoretical, based on the nature of scalar corrections. Direct empirical verification showing that AUROC and refinement remain degraded after optimal temperature scaling would strengthen this interpretation.
\end{enumerate}

\begin{enumerate}
\item \textbf{Domain-specific calibration:} Predictions P4 and P5 regarding domain-specific calibration variation and domain-conditioned temperature scaling were not tested in this iteration.
\end{enumerate}

\begin{enumerate}
\item \textbf{Causal claims:} The comparison between base and instruction-tuned models establishes association, not causation. Other differences between model versions beyond RLHF (continued pretraining, supervised fine-tuning) may contribute to the observed effects.
\end{enumerate}

\subsubsection{6.4 Implications}

For practitioners deploying language models in applications requiring reliable confidence estimates (selective prediction, human-AI collaboration, uncertainty-aware decision making), these findings suggest:

\begin{enumerate}
\item AUROC for margin-based correctness prediction may be a useful complement to ECE for evaluating confidence reliability.
\end{enumerate}

\begin{enumerate}
\item Instruction-tuned models may require different calibration approaches than base models if discriminative quality is important.
\end{enumerate}

\begin{enumerate}
\item Monitoring both calibration (ECE, reliability) and discrimination (AUROC, refinement) metrics may provide a more complete picture of confidence signal quality.
\end{enumerate}
